\needspace{32mm}
\subsection*{S01}
\path{TFPT_Fortsetzung.md}\par
Korrigierte Fortsetzung: T-Faktorisierung, Vermittler, Zweizellenzweig, Uhren, Chiralität.
{\footnotesize SHA-256: \texttt{07b941c4c325a887e09dcfc88216f44a}\newline\texttt{d0603dd1b24ea1be2310995f20d005b7}\par}
\needspace{32mm}
\subsection*{S02}
\path{TFPT_Omega_Praeparation_Mehrzeittest_2026-09-14.md}\par
Explizite Präparation, Aufzeichnung, Rohwahrscheinlichkeiten, Ressourcen und Mehrzeittest.
{\footnotesize SHA-256: \texttt{065fa53688b8518f81e8d9ea9e9e5a55}\newline\texttt{7e44a597456bf2b29a49b9db2d6bc724}\par}
\needspace{32mm}
\subsection*{S03}
\path{TFPT_TOE_GEMEINSAMER_URSPRUNG_2026-09-14.md}\par
Gemeinsamer Ursprung: Zweizellenenergie, Verschränkung, Uhrsignal, Propagatorbrücke und TOE-Vertrag.
{\footnotesize SHA-256: \texttt{864bd44fa12d8deb88cb359742ee908b}\newline\texttt{8b5cc8a12cd2d39bc22c6fdced757b2e}\par}
\needspace{32mm}
\subsection*{S04}
\path{TFPT_UNIVERSALRAUM_GESAMTSYNTHESE_2026-09-14.md}\par
Frühe Gesamtsynthese mit Fixpunkthypothese, Zahlenübersicht und Einordnung zu RH und Komplexität; mehrere Aussagen später korrigiert.
{\footnotesize SHA-256: \texttt{f62f91af70b43ce349a9e4c17f9e4887}\newline\texttt{7099ae9b6c992f0c1482f4eb77ee6694}\par}
\needspace{32mm}
\subsection*{S05}
\path{Analyse_und_Rekonstruktion.md}\par
Gegenproben, harte und weiche Hülle, Vermittler, Viereruhr, GNS und physikalische Grenzen.
{\footnotesize SHA-256: \texttt{5726d25cb03443c93cb18d28a8b1589d}\newline\texttt{b0eee06415a84995251875bb00883d60}\par}
\needspace{32mm}
\subsection*{S06}
\path{TFPT_Universalraum_Rekonstruktion_2026-09-14.md}\par
Quotienten 45/30, Register- und Quellenprüfung, Stabilizerhindernis, E8-Rückrekonstruktion.
{\footnotesize SHA-256: \texttt{4515aadc06b13c0216994bb3036bb201}\newline\texttt{380d00128ab4ce5e6f19f1e7e230fc1e}\par}
\needspace{32mm}
\subsection*{S07}
\path{TFPT_UNIVERSALRAUM_INVERSION_2026-09-14.md}\par
Inversionshypothese, Matchingausführung, Registerrekurrenz, Clockblockade und endliche Faktorprüfung.
{\footnotesize SHA-256: \texttt{b86f846bf2b20a78d2d8997bd303fe5f}\newline\texttt{4da82497c23f969e8e616e4d9361e95b}\par}
\needspace{32mm}
\subsection*{S08}
\path{tfpt_anschluss_zellen_seam_2026-09-14.pdf}\par
Zellhülle, Tick-Lifts, Kettennumerik, Rotor/CAR-Zustandsgate, Pati-Salam-Läufe und Indexpräzisierungen.
{\footnotesize SHA-256: \texttt{6029544e7bb9ee2e3fa7f94717fa1d18}\newline\texttt{61b5d0b598e7543cd8abcde1ddf0e8c8}\par}
\needspace{32mm}
\subsection*{S09}
\path{tfpt_compiler_universalraum_2026-09-13.pdf}\par
Grundmanuskript Version 1.0: Architektur, Standardmodell, Zahlen, Prozess, Register und T1 bis T8.
{\footnotesize SHA-256: \texttt{48ea9ae625ce418607eb718478525d80}\newline\texttt{2c52844183745f27849e17dd63bad43a}\par}
\needspace{32mm}
\subsection*{S10}
Eingefügter Text: Erklärung des Universalraums\par
Gesprächstext: einfache Universalraum-Erklärung, hypothetische Anwendungen, Zuse und Wolfram.
{\footnotesize SHA-256: \texttt{5c314fd3f27634dfd0cc8e2ed56e5e64}\newline\texttt{2a457e08df9db9558ee6c962a7f8fd59}\par}
\needspace{32mm}
\subsection*{S11}
Eingefügter Text: Die fünf gesetzten Fugen\par
Fünf Fugen: Clebsch-Graph, Austausch, Sektortrennung, Orbitargument und Z4-Erweiterung.
{\footnotesize SHA-256: \texttt{79258cd3d27f5a8219ef3b7d597978ea}\newline\texttt{eb6c2c1f0fd18cb636f0a70bbd548ca8}\par}